% SOURCE_STATEMENT: "Is the Gregory–Laflamme instability the only mechanism by which general relativity changes horizon topology? In spacetime dimension $D>4$, is every stable black hole a Myers–Perry black hole satisfying $J/M^{D-3}=O(1)$?"
% Preparation and final self-audit: CONTRIBUTING.md prompts C and D; source/audit artifacts retained in the local physics_statements_2026_10_11 batch.
% !TeX program = lualatex
\documentclass[11pt]{article}
\usepackage[margin=25mm]{geometry}
\usepackage{fontspec}
\setmainfont{Latin Modern Roman}
\usepackage{amsmath,amssymb,mathtools}
\usepackage{unicode-math}
\setmathfont{Latin Modern Math}
\usepackage{xurl,hyperref}
\hypersetup{colorlinks=true,urlcolor=blue}
\setcounter{secnumdepth}{0}
\setlength{\parindent}{0pt}
\setlength{\parskip}{5pt}
\setlength{\emergencystretch}{3em}
\newcommand{\sref}[2]{\hyperlink{source:#1:#2}{[S#2]}}
\newcommand{\cataloguescope}[1]{\paragraph{Recorded target.}#1}
\newcommand{\R}{\mathbb R}
\newcommand{\C}{\mathbb C}
\newcommand{\N}{\mathbb N}
\newcommand{\Z}{\mathbb Z}
\newcommand{\Q}{\mathbb Q}
\newcommand{\F}{\mathbb F}
\newcommand{\eps}{\varepsilon}
\begin{document}
\section{Horizon Topology Changes and Stable Higher-Dimensional Black Holes}\label{statement}

\subsection{Definitions and mathematical statement}

\paragraph{Shared notation.}$\mathbb N=\{1,2,\ldots\}$; the standard number fields have their usual meanings. All additional conventions are stated locally.


\paragraph{Definitions and conventions.}
The Gregory--Laflamme instability is an instability of extended higher-dimensional black objects associated with long-wavelength perturbations. Horizon topology is the topology of a horizon cross-section; a topology-changing evolution must distinguish apparent and event horizons. Myers--Perry metrics are the higher-dimensional rotating vacuum black-hole family. Stability requires a chosen perturbation class. In the source formula, $M$, $J$ and $D$ denote mass, angular momentum and spacetime dimension, and $O(1)$ denotes a bound in its chosen scaling convention.

\paragraph{Target retained from the canonical source catalogue.}
\paragraph{Catalogue wording retained for traceability.}
Is the Gregory–Laflamme instability the only mechanism by which general relativity changes horizon topology? In spacetime dimension $D>4$, is every stable black hole a Myers–Perry black hole satisfying $J/M^{D-3}=O(1)$?

\paragraph{Source scope and qualification.}
Both questions are retained as source conjectural speculation. The supplied $J/M^{D-3}$ is not replaced by a guessed dimensionless normalization; conventional dimensionless spin uses dimension-dependent powers. The asymptotic and rotation hypotheses need source clarification.
\subsection{Short English statement}
Are there other mechanisms for horizon topology changes, and are all stable higher-dimensional black holes moderately rotating Myers--Perry solutions?
\subsection{Sources}
\begin{itemize}
\item[\hypertarget{source:statement:1}{S1}] ULAM.AI and the UnsolvedMath Contributors, UnsolvedMath, catalogue entry. Catalogue curation under CC BY 4.0; definitions and English exposition added locally.
\url{https://huggingface.co/datasets/ulamai/UnsolvedMath}.
\item[\hypertarget{source:statement:2}{S2}] Mathematical General Relativity (2018), Oberwolfach Reports, DOI 10.4171/owr/2018/36, pp. 2157--2251. Primary report accessed 2026-10-11; relevant PDF pages 37.
\url{https://doi.org/10.4171/owr/2018/36}.
\item[\hypertarget{source:statement:3}{S3}] Emparan--Reall higher-dimensional black holes (catalogue-listed context; not independently consulted in this preparation).
\url{https://doi.org/10.1007/978-3-642-41992-8}.
\item[\hypertarget{source:statement:4}{S4}] Black holes galore in D greater than 4 (catalogue-listed context; not independently consulted in this preparation).
\url{https://doi.org/10.1007/s10714-025-03398-x}.
\end{itemize}
\label{end:statement}
\end{document}
